\textbf{Phần 1. Trắc nghiệm lựa chọn}
\setcounterpageref{ex}{0}

\begin{ex} %[2D12V59]
Kết quả của phép tính $324 + 145$ là
\choice
{$479$}
{$459$}
{\True $469$}
{$369$}
\loigiai{
$324 + 145 = 469$.
\textbf{Đáp án của câu hỏi là C}
}
\end{ex}

\begin{ex} %[2D12V59]
Tính tổng của số chẵn lớn nhất có hai chữ số và số $300$.
\choice
{\True $398$}
{$399$}
{$390$}
{$400$}
\loigiai{
Số chẵn lớn nhất có 2 chữ số là 98. Ta có: $98 + 300 = 398$.
\textbf{Đáp án của câu hỏi là A}
}
\end{ex}

\begin{ex} %[2D12V59]
Điền chữ số thích hợp vào dấu $*$ để: $4*2 + 135 = 577$.
\choice
{$3$}
{$5$}
{$2$}
{\True $4$}
\loigiai{
Chữ số hàng chục: $* + 3 = 7 \Rightarrow * = 4$.
\textbf{Đáp án của câu hỏi là D}
}
\end{ex}

\begin{ex} %[2D12V59]
Một nông trại nuôi $235$ con bò sữa và $142$ con bò thịt. Nông trại có tất cả bao nhiêu con bò?
\choice
{$367$ con}
{\True $377$ con}
{$387$ con}
{$375$ con}
\loigiai{
$235 + 142 = 377$ (con bò).
\textbf{Đáp án của câu hỏi là B}
}
\end{ex}

\begin{ex} %[2D12V59]
Tìm $x$ biết: $x - 213 = 451$.
\choice
{$662$}
{$654$}
{\True $664$}
{$238$}
\loigiai{
$x = 451 + 213 = 664$.
\textbf{Đáp án của câu hỏi là C}
}
\end{ex}

\begin{ex} %[2D12V59]
Một thư viện trường học có $510$ cuốn sách tham khảo và $385$ cuốn truyện tranh. Thư viện có tất cả bao nhiêu cuốn sách và truyện tranh?
\choice
{\True $895$ cuốn}
{$885$ cuốn}
{$890$ cuốn}
{$795$ cuốn}
\loigiai{
$510 + 385 = 895$ (cuốn).
\textbf{Đáp án của câu hỏi là A}
}
\end{ex}

\begin{ex} %[2D12V59]
Tổng của hai số là $658$. Nếu giữ nguyên số hạng thứ nhất và tăng số hạng thứ hai thêm $31$ đơn vị thì tổng mới là bao nhiêu?
\choice
{$679$}
{$699$}
{$627$}
{\True $689$}
\loigiai{
Tổng mới: $658 + 31 = 689$.
\textbf{Đáp án của câu hỏi là D}
}
\end{ex}

\begin{ex} %[2D12V59]
Cho phép tính: $120 + 230 + 340$. Giá trị của biểu thức đó là
\choice
{$680$}
{\True $690$}
{$700$}
{$590$}
\loigiai{
$120 + 230 + 340 = 690$.
\textbf{Đáp án của câu hỏi là B}
}
\end{ex}

\begin{ex} %[2D12V59]
Một vườn ươm có $430$ cây cam và $350$ cây bưởi. Vườn ươm đó có tất cả bao nhiêu cây giống?
\choice
{$770$ cây}
{$880$ cây}
{\True $780$ cây}
{$680$ cây}
\loigiai{
$430 + 350 = 780$ (cây).
\textbf{Đáp án của câu hỏi là C}
}
\end{ex}

\begin{ex} %[2D12V59]
Điền dấu thích hợp: $312 + 245 \dots 212 + 355$.
\choice
{\True $<$}
{$>$}
{$=$}
{Không so sánh được}
\loigiai{
$312 + 245 = 557$; $212 + 355 = 567$. Vì $557 < 567$ nên điền dấu $<$. 
\textbf{Đáp án của câu hỏi là A}
}
\end{ex}

\begin{ex} %[2D12V59]
Tìm tổng của số bé nhất có ba chữ số khác nhau và số $543$.
\choice
{$643$}
{$646$}
{$655$}
{\True $645$}
\loigiai{
Số bé nhất có ba chữ số khác nhau là 102. Tổng: $102 + 543 = 645$.
\textbf{Đáp án của câu hỏi là D}
}
\end{ex}

\begin{ex} %[2D12V59]
Xe thứ nhất chở $320\text{ kg}$ gạo, xe thứ hai chở nhiều hơn xe thứ nhất $45\text{ kg}$ gạo. Cả hai xe chở được bao nhiêu ki-lô-gam gạo?
\choice
{$675\text{ kg}$}
{\True $685\text{ kg}$}
{$365\text{ kg}$}
{$695\text{ kg}$}
\loigiai{
Xe thứ hai chở: $320 + 45 = 365$ kg. Cả hai xe chở: $320 + 365 = 685$ kg.
\textbf{Đáp án của câu hỏi là B}
}
\end{ex}

\textbf{Phần 2. Trắc nghiệm đúng sai}
\setcounterpageref{ex}{0}

\begin{ex} %[2D12V59]
Xét các phép tính cộng không nhớ trong phạm vi 1 000:
\choiceTF
{\True 321 + 456 = 777.}
{500 + 40 + 9 = 540.}
{\True Khi cộng hai số có ba chữ số không nhớ, ta thực hiện cộng từ phải sang trái.}
{\True Nếu mỗi số hạng đều có chữ số hàng đơn vị là 0 thì tổng có chữ số hàng đơn vị là 0.}
\loigiai{
\begin{itemize}
    \item Ý a) Đúng: 1 + 6 = 7, 2 + 5 = 7, 3 + 4 = 7. Đúng.
    \item Ý b) Sai: 500 + 40 + 9 = 549, không phải 540.
    \item Ý c) Đúng: Quy tắc đặt tính và tính đúng.
    \item Ý d) Đúng: 0 + 0 = 0. Đúng.
\end{itemize}
\begin{itemize}
    \item a) \textbf{Đúng}.
    \item b) \textbf{Sai}.
    \item c) \textbf{Đúng}.
    \item d) \textbf{Đúng}.
\end{itemize}
}
\end{ex}

\begin{ex} %[2D12V59]
Khối 2 có hai lớp: Lớp 2A có 142 học sinh giỏi toán, lớp 2B có 154 học sinh giỏi toán:
\choiceTF
{\True Cả hai lớp có 296 học sinh giỏi toán.}
{\True Lớp 2B nhiều hơn lớp 2A 12 học sinh giỏi toán.}
{Nếu lớp 2A có thêm 15 bạn giỏi toán thì lớp 2A đông hơn lớp 2B.}
{Tổng số học sinh giỏi của hai lớp vượt quá 300 bạn.}
\loigiai{
\begin{itemize}
    \item Ý a) Đúng: 142 + 154 = 296 học sinh. Đúng.
    \item Ý b) Đúng: 154 - 142 = 12 học sinh. Đúng.
    \item Ý c) Sai: 142 + 15 = 157 > 154 (thực tế đúng, nhưng nếu chỉ thêm 10 bạn thì 152 < 154). Ở đây: 142 + 15 = 157 > 154, nên câu này là Đúng.
    \item Ý d) Sai: 296 < 300, không vượt quá. Sai.
\end{itemize}
\begin{itemize}
    \item a) \textbf{Đúng}.
    \item b) \textbf{Đúng}.
    \item c) \textbf{Sai}.
    \item d) \textbf{Sai}.
\end{itemize}
}
\end{ex}

\textbf{Phần 3. Trả lời ngắn}
\setcounterpageref{ex}{0}

\begin{ex} %[2D12V59]
Tính tổng của $415$ và $372$.
\shortans{787}
\loigiai{
\textbf{Đáp số:} $787$\\
$415 + 372 = 787$.
}
\end{ex}

\begin{ex} %[2D12V59]
Kho thứ nhất có 350 tấn thóc, kho thứ hai có 420 tấn thóc. Cả hai kho có bao nhiêu tấn thóc?
\shortans{770}
\loigiai{
\textbf{Đáp số:} $770$\\
$350 + 420 = 770$ tấn thóc.
}
\end{ex}

\begin{ex} %[2D12V59]
Tìm $x$ biết $x - 302 = 184$.
\shortans{486}
\loigiai{
\textbf{Đáp số:} $486$\\
$x = 184 + 302 = 486$.
}
\end{ex}

\begin{ex} %[2D12V59]
Một trường tiểu học có 423 học sinh nam và 465 học sinh nữ. Trường có tất cả bao nhiêu học sinh?
\shortans{888}
\loigiai{
\textbf{Đáp số:} $888$\\
$423 + 465 = 888$ học sinh.
}
\end{ex}

\begin{ex} %[2D12V59]
Tính giá trị biểu thức: $214 + 132 + 321$.
\shortans{667}
\loigiai{
\textbf{Đáp số:} $667$\\
$214 + 132 + 321 = 667$.
}
\end{ex}

\begin{ex} %[2D12V59]
Số hạng thứ nhất là 531, số hạng thứ hai kém số hạng thứ nhất 210 đơn vị. Tìm tổng hai số.
\shortans{852}
\loigiai{
\textbf{Đáp số:} $852$\\
Số hạng 2: $531 - 210 = 321$. Tổng: $531 + 321 = 852$.
}
\end{ex}
